\subsubsection{Torsion equality cases and the explicit fierce cover}

\begin{lemma}[Proved: killing a degree-zero obstruction]
\label{D:torsion-first-obstruction-cover}
Let $C/k$ be a smooth proper connected curve over an
algebraically closed field of characteristic $p>0$, let
$L$ be a torsion line bundle of degree zero, and let
$e\in H^1(C,L)$.  There is a finite map $h:C'\to C$,
obtained by a sequence of finite \'etale covers and relative
Frobenius morphisms, such that $h^*L\simeq\mathcal O_{C'}$
and $h^*e=0$.  If the product of the \'etale degrees is $s$
and the product of the Frobenius degrees is $q$, then
\[
             \deg h=qs,\qquad g(C')-1=s(g(C)-1).
\]
The Frobenius degree can be increased arbitrarily after
these two vanishing properties have been achieved.
\end{lemma}

\begin{proof}
Write the order of $L$ as $ap^b$, with $(a,p)=1$.
Pull back first by relative Frobenius of degree $p^b$;
with the necessary $k$-twists, this raises the line bundle
to its $p^b$-th power.  Its order now divides $a$.
A trivialization of its $a$-th power defines a finite
\'etale $\mu_a$-torsor on which it trivializes, by
\cite[Lemmas 59.28.1 and 59.28.5, section 03PK]
{StacksBFiniteCone}.  Choose a connected component of this
torsor.  We have reduced the assertion to killing one class
in $H^1(D,\mathcal O_D)$ for a smooth proper connected curve
$D$.

Let $F$ denote the semilinear Frobenius on this finite
dimensional space.  Since $k$ is perfect, the images of
$F^r$ are $k$-subspaces.  For sufficiently large $r$ the
images stabilize, and the space is the direct sum of a
nilpotent part and a part on which $F$ is bijective.
Pullback by an appropriate relative Frobenius kills the
nilpotent part of our class.  We explain why the remaining
part is a $k$-linear combination of Frobenius-fixed classes.

On an $r$-dimensional bijective semilinear space, write
$F(x)=A x^{[p]}$ with $A$ invertible.  The polynomial map
$F-1$ is \'etale, with Jacobian $-I$.  Its kernel has
$p^r$ points: after multiplying the equations by $A^{-1}$
they are $x_i^p$ minus linear forms, whose relatively
prime leading monomials give the basis
$\prod x_i^{e_i}$, $0\leq e_i<p$, for the quotient algebra.
This is also the exact calculation in
\cite[Lemma 59.63.2, section 0A3J]{StacksBFiniteCone}.
An $\mathbb F_p$-basis of this kernel is $k$-linearly
independent.  Indeed a shortest normalized $k$-linear
relation, after applying $F$ and subtracting, would give
a shorter one unless all its coefficients lie in
$\mathbb F_p$, contradicting independence over that field.
The fixed vectors therefore form a $k$-basis of the
bijective part.

By the Artin--Schreier exact sequence, each such fixed
class in $H^1(D,\mathcal O_D)$ is the image of a
$\mathbb Z/p\mathbb Z$-torsor on $D$; see
\cite[section 59.63, tag 0A3J]{StacksBFiniteCone}.
Pullback to that finite \'etale torsor kills its class.
A connected component of the fibre product of the
finitely many torsors kills every basis class and hence
the chosen $k$-linear combination.  This proves the
vanishing assertion.  All Frobenius maps here mean
$k$-linear relative Frobenius maps from the inverse
Frobenius twist; the usual semilinear notation records
the induced action under the corresponding scalar twists.
Relative Frobenius preserves the genus, and an \'etale
map of degree $s_i$ multiplies $g-1$ by $s_i$.
Multiplying these identities proves the displayed
formulas.  Further Frobenius pullbacks preserve a zero
class and a trivial line bundle.
\end{proof}

\begin{corollary}[Proved: the torsion equality case]
\label{D:negative-curve-equality-finite-wild}
In the surface setting of
Corollary~\ref{D:negative-curve-finite-wild}, suppose $g\geq2$,
$d=2g-2$, and $\omega_C\otimes N_{C/Y}$ is torsion.
Then the full exceptional wild inertia has finite image.
In particular this holds when the surface germ is
$\mathbb Q$-Gorenstein, without any restriction on its
Cartier index.
\end{corollary}

\begin{proof}
The line bundle
$L=T_C\otimes N^{-1}=(\omega_C\otimes N)^{-1}$
has degree zero and is torsion.  Apply
Lemma~\ref{D:torsion-first-obstruction-cover} to the
first-neighborhood class $e_Y$.  Write the resulting
degree as $qs$ as there, and add one relative Frobenius
if necessary so that $q>1$.
For $N'=h^*N$ and $n\geq2$, both line bundles in
\eqref{D:cone-tangent}--\eqref{D:cone-normal} have degree
\[
                 qs(n-1)d.
\]
This exceeds $sd=2g(C')-2$, since $q>1$.
Their first cohomology vanishes, and
\eqref{D:cone-first} holds by construction.
The finite cone cover and finite-wild propositions apply.

For the last assertion write
$K_Y=f^*K_S+aC$.  Adjunction gives
$a=-1-(2g-2)/d=-2$.
If $rK_S$ is Cartier, it is trivial on the local scheme
$S$.  Consequently $\mathcal O_Y(r(K_Y+2C))$ is
trivial, and restricting to $C$ gives
\[
                    (\omega_C\otimes N)^{\otimes r}
                         \simeq\mathcal O_C.
\]
This proves the torsion hypothesis directly; no cyclic
canonical cover is required.
\end{proof}

\begin{corollary}[Proved: full wild bounds in the explicit genus-two example]
\label{D:wild-genus2-finite-inertia}
For the explicit surfaces $T\to S$ in
Example~\ref{D:wild-genus2-single-curve}, the image of
exceptional wild inertia in the full punctured
fundamental group has order at most four for $T$ and
at most eight for $S$.  These statements also hold for
their punctured strict henselizations.
\end{corollary}

\begin{proof}
The source resolution in
Proposition~\ref{D:wild-genus2-source} has exceptional curve
\[
                   C':\quad w^2+w=z^5,
\]
of genus two and square $-2$.  Let $P_0,P_1$ be its two
points with $z=0$, and let $\infty$ be its point at infinity.
The already checked divisor of the function $c$ gives
\[
                    N_{C'/V}\simeq\mathcal O_{C'}(-P_0-P_1).
\]
On this curve
$\operatorname{div}(z)=P_0+P_1-2\infty$.
Also $\operatorname{div}(dz)=2\infty$.  For completeness,
on the infinity chart $u=z^{-1}$, $v=w/z^3$, the equation
is $v^2+u^3v=u$, and differentiation in characteristic two
gives
\[
                  (1+u^2v)\,du=u^3\,dv.
\]
Here $v$ is a uniformizer, $\operatorname{ord}_\infty u=2$,
and hence $\operatorname{ord}_\infty dz=2$; on the finite
chart $z$ is an \'etale coordinate, so there are no other
zeros or poles.  It follows that
$N_{C'/V}\simeq\omega_{C'}^{-1}$ and
$T_{C'}\otimes N_{C'/V}^{-1}\simeq\mathcal O_{C'}$.

The two-chart calculation in
Corollary~\ref{D:wild-genus2-pro2} gives
$F(w/z)=w/z^2$ and $F(w/z^2)=0$ on
$H^1(C',\mathcal O_{C'})$.  Thus the first-neighborhood
class is killed by the relative Frobenius of degree four.
For that map, the two later line-bundle degrees are both
$8(n-1)>2$ for $n\geq2$.  Lemma~\ref{D:finite-formal-cone-cover}
therefore supplies a fixed cone cover of degree four
over the completed germ $T$.
Proposition~\ref{D:finite-cone-wild-bound} gives the first
bound.

Composing that finite local cone algebra with the checked
quasi-\'etale double cover $T\to S$ gives a fixed finite
local cone cover of degree eight over the completed germ
$S$.  Its actual radial DVR extends the exceptional
DVR of $S$, by the finite morphism of resolutions in
Proposition~\ref{D:wild-genus2-source}.  Every punctured
cover of $S$ again pulls back to a punctured cone cover,
so the identical inertia-index argument gives the
second bound.  The first example already has a nonzero
wild image of order two; the present upper bounds do
not assert that four or eight is attained.
\end{proof}

\begin{remark}[Exact companion certificate]
The program \nolinkurl{parts/D/code/check_finite_cone_genus2.py}
checks sixteen characteristic-two polynomial, divisor-degree,
Frobenius-matrix and obstruction-degree assertions used in
Corollary~\ref{D:wild-genus2-finite-inertia}.
All passed in the independent D rerun in
$0.0381131569884019$ seconds; the output is
\nolinkurl{parts/D/output/finite_cone_genus2.json}.
Root independently read and reran the same program, with all
sixteen assertions passing in $0.04670357100258116$ seconds.
The formal-cover construction and inertia conclusions are
proved separately above.
\end{remark}

